% ============================================================================
% CHAPITRE 8 : MISE EN ÉQUATIONS ET SYSTÈMES ALGÉBRIQUES (PROBLÈMES TEXTUELS)
% ============================================================================
\chapter{Mise en Équations et Systèmes Algébriques (Problèmes Textuels)}

La \textbf{Question~3} de l'épreuve d'Algèbre à l'examen d'admission de l'UCLouvain est une grande question de \textbf{modélisation mathématique} : un récit concret d'une demi-page à une page décrit une situation mettant souvent en scène \emph{Charles, Laurent, Aude, Sophie} ou un projet de l'EPL. Il s'agit de traduire méthodiquement chaque phrase de l'énoncé en une équation algébrique, puis de résoudre le système obtenu (linéaire, ensembliste ou bilinéaire) sans calculatrice.

Ce chapitre expose la méthode de décodage en 5 étapes et analyse en détail \textbf{10 problèmes complets de Question~3 posés entre 2022 et 2026}.

% ----------------------------------------------------------------------------
\section{Méthodologie en 5 étapes pour décoder un problème textuel}
% ----------------------------------------------------------------------------

Face à un long texte, l'erreur la plus fréquente consiste à essayer de tout garder en tête ou à écrire des calculs numériques épars sans définir ses variables. Appliquez toujours ce protocole d'ingénieur :

\begin{theoreme}[Protocole en 5 étapes pour réussir la Question 3]
\begin{enumerate}
    \item \textbf{Étape 1 -- Définition explicite des inconnues (et de leurs unités !) :}
    Dès les premières lignes de votre copie, déclarez chaque variable par une phrase claire en précisant son \textbf{unité} (kg, €, m, $\text{m}^3$, nombre entier de sacs ou de sièges).
    \emph{Conseil :} Choisissez des lettres parlantes ($F, B, V$ pour les ballons de \textbf{F}oot, \textbf{B}asket, \textbf{V}olley ; $p, d, b$ pour \textbf{p}lat, \textbf{d}essert, \textbf{b}oisson) plutôt que $x_1, x_2, x_3, x_4$ : vous éviterez ainsi de confondre vos variables en fin de calcul !
    \item \textbf{Étape 2 -- Traduction phrase par phrase en équations numérotées $(1), (2), (3)\dots$ :}
    Surlignez chaque donnée chiffrée ou relation de proportionnalité dans le texte (« trois fois plus que », « un sixième de moins », « la moyenne de », « un nombre impair dont la majorité est $\frac{N+1}{2}$ ») et écrivez l'équation correspondante.
    \begin{itemize}
        \item \textbf{Homogénéité des unités :} Vérifiez que les deux membres de chaque équation sont exprimés dans la \textbf{même unité} (ne mélangez pas des grammes avec des kilogrammes, ou des centilitres avec des litres : convertissez $30\text{ kg} = 30\,000\text{ g}$ et $94 \times 25\text{ cL} = 2350\text{ cL}$ !).
    \end{itemize}
    \item \textbf{Étape 3 -- Simplification immédiate et choix de la variable pivot :}
    Avant de vous lancer dans de grandes substitutions, divisez chaque équation par son plus grand facteur commun (ex. $4p + 4d + 4b = 120 \implies p + d + b = 30$). Repérez si plusieurs inconnues s'expriment simplement en fonction d'une seule \textbf{variable pivot} (par exemple $E = 3C$, $D = 2B$, ou $B = 2V$ et $F = \frac{3}{2}V$).
    \item \textbf{Étape 4 -- Résolution par substitution en cascade ou combinaisons linéaires :}
    \begin{itemize}
        \item Pour un \textbf{système linéaire} ($ax + by + cz = d$), procédez par \textbf{combinaisons linéaires} (soustraire deux équations ayant des termes identiques comme $(1) - (2)$ pour éliminer $c_0$ ou $p + d$) ou par substitution de la variable pivot.
        \item Pour un \textbf{système bilinéaire} (comportant des produits d'inconnues comme $N \cdot p = 3240$ et $(N + 5)p = 3510$), développez $(N + 5)p = Np + 5p$ et remplacez le bloc $Np$ tout entier par sa valeur $3240$ !
    \end{itemize}
    \item \textbf{Étape 5 -- Contrôle final dans l'énoncé et report dans les cadres :}
    En 30 secondes, vérifiez que vos valeurs numériques redonnent bien les totaux de l'énoncé, puis inscrivez lisiblement vos réponses (avec unités) dans les cadres prévus.
\end{enumerate}
\end{theoreme}

% ----------------------------------------------------------------------------
\section{Systèmes linéaires classiques et variable pivot (2022, 2023, 2024, 2025, 2026)}
% ----------------------------------------------------------------------------

\subsection{Élimination rapide par soustraction d'équations}

\begin{exemple}[L'addition au restaurant à la Mer du Nord -- UCLouvain, Septembre 2022, Q3]
\textbf{Résumé de l'énoncé :} Au restaurant, tous les plats principaux sont au même prix $p$, tous les desserts au prix $d$ et toutes les boissons au prix $b$.
\begin{itemize}
    \item Une affiche annonce : « Formule familiale : plat + boisson + dessert pour $4$ personnes à $100$\,€ au lieu de $120$\,€ ! » (le prix normal pour $4$ formules complètes est donc de $120$\,€).
    \item Aude a pris $1$ plat et $2$ boissons. Laurent a pris $1$ plat, $1$ boisson et $1$ dessert. Charles a pris autant de boissons qu'Aude et Laurent ensemble (soit $2 + 1 = 3$ boissons) et $1$ dessert, sans plat.
    \item Charles a réglé l'addition totale de $80$\,€ pour les trois amis.
    \item Aude a remarqué : « Si on prend tous les trois un dessert, cela revient aussi cher que ton plat et le mien ! »
\end{itemize}
Combien chacun doit-il payer (et combien Aude et Laurent doivent-ils rembourser à Charles) ?

\textbf{Résolution structurée :}
\begin{enumerate}
    \item \textbf{Inconnues :} $p$ = prix d'un plat (€), $d$ = prix d'un dessert (€), $b$ = prix d'une boisson (€).\\
    Parts respectives : Laurent $= p + b + d$, Aude $= p + 2b$, Charles $= 3b + d$.
    \item \textbf{Mise en équations :}
    \begin{align}
        \text{Prix normal (4 formules) :}\quad 4p + 4d + 4b = 120 &\iff p + d + b = 30 \label{eq:rest1}\\
        \text{Facture totale (3 amis) :}\quad 2p + 2d + 6b = 80 &\iff p + d + 3b = 40 \label{eq:rest2}\\
        \text{Remarque d'Aude (3 desserts = 2 plats) :}\quad 3d &= 2p \label{eq:rest3}
    \end{align}
    \item \textbf{Résolution éclair :}
    \begin{itemize}
        \item Notons d'emblée que l'équation~\eqref{eq:rest1} $p + d + b = 30$ donne \textbf{directement la part de Laurent : $30$\,€} !
        \item En soustrayant \eqref{eq:rest2} $-$ \eqref{eq:rest1} : $(p + d + 3b) - (p + d + b) = 40 - 30 \iff 2b = 10 \iff \mathbf{b = 5\,\text{€}}$.
        \item On remplace $b = 5$ dans \eqref{eq:rest1} : $p + d = 25$. Avec \eqref{eq:rest3} ($p = \frac{3}{2}d$), il vient $\frac{5}{2}d = 25 \iff \mathbf{d = 10\,\text{€}}$ et $\mathbf{p = 15\,\text{€}}$.
    \end{itemize}
    \item \textbf{Conclusion :} Laurent doit payer $p + b + d = 15 + 5 + 10 = \boxed{30\,\text{€}}$, Aude doit payer $p + 2b = 15 + 10 = \boxed{25\,\text{€}}$, et Charles a consommé pour $3b + d = 15 + 10 = \boxed{25\,\text{€}}$ (ayant payé les $80$\,€, il doit recevoir $30$\,€ de Laurent et $25$\,€ d'Aude).
\end{enumerate}
\end{exemple}

\begin{exemple}[Logistique de brasserie : bouteilles, casiers, fûts et palettes -- UCLouvain, Juillet 2023 Série 1, Q3]
\textbf{Résumé de l'énoncé :} Chez \emph{Eaux, Pils et Liqueurs (EPL)}, la bière ($1\text{ L} = 1\text{ kg}$) est vendue en casiers de $24$ bouteilles de $25\text{ cL}$ (soit $24 \times 0,25 = 6\text{ L} = 6\text{ kg}$ de bière par casier plein) ou en fûts de $50\text{ L}$ (soit $50\text{ kg}$ de bière par fût plein).
\begin{itemize}
    \item Charles rapporte $6$ casiers de bouteilles vides et repart avec $8$ casiers pleins : sa voiture pèse $66\text{ kg}$ de plus à la sortie.
    \item Laurent se fait livrer $1$ palette de $8$ fûts pleins et $2$ palettes de $60$ casiers pleins chacune : le poids total (incluant les $3$ palettes) est de $2\,370\text{ kg}$.
    \item Le poids d'un fût plein équivaut à celui de $4$ casiers pleins, et le poids d'un fût vide équivaut à celui de $40$ bouteilles vides.
\end{itemize}
Déterminer le poids d'une bouteille vide ($b$), d'un casier sans bouteilles ($c$), d'un fût vide ($f$) et d'une palette seule ($p$).

\textbf{Résolution par variables agrégées (l'astuce qui simplifie tout !) :}
Plutôt que d'écrire tout de suite des équations compliquées en $b, c, f, p$, introduisons d'abord les poids intermédiaires d'un \textbf{casier de vidanges} $C_{\text{vide}} = c + 24b$, d'un \textbf{casier plein} $C_{\text{plein}} = C_{\text{vide}} + 6$ (car il contient $6\text{ kg}$ de bière) et d'un \textbf{fût plein} $F_{\text{plein}} = f + 50$ :
\begin{enumerate}
    \item \textbf{Voiture de Charles :} $8\,C_{\text{plein}} - 6\,C_{\text{vide}} = 66 \iff 8(C_{\text{vide}} + 6) - 6\,C_{\text{vide}} = 66 \iff 2\,C_{\text{vide}} + 48 = 66 \iff \mathbf{C_{\text{vide}} = 9\text{ kg}}$, et donc $\mathbf{C_{\text{plein}} = 9 + 6 = 15\text{ kg}}$.
    \item \textbf{Fût plein et fût vide :} $F_{\text{plein}} = 4\,C_{\text{plein}} = 4 \times 15 = \mathbf{60\text{ kg}}$. Comme il contient $50\text{ kg}$ de bière, le fût vide pèse $\mathbf{f = 60 - 50 = 10\text{ kg}}$.
    \item \textbf{Bouteille vide et casier vide :} $f = 40b \iff 40b = 10 \iff \mathbf{b = 0,25\text{ kg}}$. Puis de $C_{\text{vide}} = c + 24b = 9$, on tire $c = 9 - 24(0,25) = 9 - 6 = \mathbf{3\text{ kg}}$.
    \item \textbf{Camion de Laurent ($3$ palettes + $8$ fûts pleins + $120$ casiers pleins) :}
    \[
        3p + 8\,F_{\text{plein}} + 120\,C_{\text{plein}} = 2370 \iff 3p + 480 + 1800 = 2370 \iff 3p = 90 \iff \mathbf{p = 30\text{ kg}}.
    \]
\end{enumerate}
Réponses : $\boxed{b = 0,25\text{ kg},\quad c = 3\text{ kg},\quad f = 10\text{ kg},\quad p = 30\text{ kg}}$.
\end{exemple}

\subsection{Utilisation d'une variable pivot dans les systèmes à 4, 5 ou 6 inconnues}

\begin{exemple}[Les coalitions électorales à l'EPL -- UCLouvain, Juillet 2024, Q3]
\textbf{Résumé de l'énoncé :} Cinq partis ($A$ = Arch, $B$ = Bio, $C$ = Chimie, $D$ = Data, $E$ = EPL) se partagent un nombre total impair $N = A + B + C + D + E$ de sièges. Pour avoir la majorité absolue, il faut au moins $M = \frac{N + 1}{2}$ sièges (le plus petit entier strictement supérieur à $N/2$).
\begin{itemize}
    \item La coalition Data + EPL dispose de $6$ sièges de plus que la majorité : $D + E = M + 6$.
    \item La coalition Bio + Chimie + Data est tout juste à la majorité : $B + C + D = M$.
    \item EPL a trois fois plus de sièges que Chimie ($E = 3C$), Data en a deux fois plus que Bio ($D = 2B$), et Bio + Data ensemble représentent les $\frac{2}{5}$ du total des sièges ($B + D = \frac{2}{5}N$).
\end{itemize}
Combien de sièges chaque parti a-t-il obtenus ?

\textbf{Résolution par variable pivot (exprimer tout en fonction de $B$) :}
\begin{enumerate}
    \item \textbf{Exprimer $D$ et $N$ puis $M$ en fonction de $B$ :}
    Puisque $D = 2B$, on a $B + D = 3B$. Or $B + D = \frac{2}{5}N$, donc :
    \[
        3B = \frac{2}{5}N \iff \mathbf{N = \frac{15}{2}B} \qquad \text{et} \qquad \mathbf{M = \frac{N + 1}{2} = \frac{15B + 2}{4}}.
    \]
    \item \textbf{Trouver $C$ et $E$ en soustrayant les deux équations de coalition :}
    On sait que $D + E = M + 6$ et $B + C + D = M$. En soustrayant la seconde à la première ($D$ et $M$ s'éliminent !) :
    \[
        (D + E) - (B + C + D) = (M + 6) - M \iff E - B - C = 6.
    \]
    Comme $E = 3C$, cela donne $3C - B - C = 6 \iff 2C = B + 6 \iff \mathbf{C = \frac{B + 6}{2}}$.
    \item \textbf{Calculer $B$ via l'équation $B + C + D = M$ :}
    Remplaçons $D = 2B$, $C = \frac{B+6}{2}$ et $M = \frac{15B+2}{4}$ dans $B + D + C = M$ :
    \[
        3B + \frac{B + 6}{2} = \frac{15B + 2}{4} \xrightarrow{\times 4} 12B + 2B + 12 = 15B + 2 \iff \mathbf{B = 10}.
    \]
    \item \textbf{Déduction en cascade de toutes les autres inconnues :}
    $D = 2B = \mathbf{20}$, $C = \frac{10 + 6}{2} = \mathbf{8}$, $E = 3C = \mathbf{24}$, $N = \frac{15}{2}(10) = \mathbf{75}$ (et $M = \frac{75+1}{2} = 38$), puis enfin $A = N - (B + C + D + E) = 75 - (10 + 8 + 20 + 24) = 75 - 62 = \mathbf{13}$.
\end{enumerate}
Réponses : $\boxed{A = 13,\quad B = 10,\quad C = 8,\quad D = 20,\quad E = 24 \quad (\text{total } N = 75,\text{ majorité } M = 38)}$.
\end{exemple}

\begin{exemple}[Trois autres problèmes de systèmes linéaires et flux (Juillet 2025, Septembre 2025, Septembre 2026)]
\begin{enumerate}
    \item \textbf{Juillet 2025, Q3 (Le stand de jus d'oranges aux Portes Ouvertes) :}
    \emph{$94$ visiteurs boivent chacun un verre de $25\text{ cL}$ (volume total $= 94 \times 25 = 2\,350\text{ cL}$), pressé à partir de $185$ oranges pesant au total $30\text{ kg} = 30\,000\text{ g}$. Il y a trois variétés : Valencia ($x$ oranges à $150\text{ g}$ et $15\text{ cL}$/fruit), Navel ($y$ oranges à $200\text{ g}$ et $10\text{ cL}$/fruit) et Sanguines ($z$ oranges à $120\text{ g}$ et $10\text{ cL}$/fruit). Chaque sac d'oranges a le même poids $P$ (en kg), ne contient qu'une variété, et Laurent a utilisé $1$ sac de Valencia de plus que de Navel.}\\
    \textbf{Étape 1 (Système $3 \times 3$ pour le nombre d'oranges $x, y, z$) :}
    \[
        \begin{cases}
            x + y + z = 185 & (1) \text{ [Nombre d'oranges]} \\
            15x + 10y + 10z = 2350 \iff 3x + 2(y + z) = 470 & (2) \text{ [Volume de jus en cL]} \\
            150x + 200y + 120z = 30\,000 \iff 15x + 20y + 12z = 3000 & (3) \text{ [Masse totale en g]}
        \end{cases}
    \]
    De $(1)$, $y + z = 185 - x$. En injectant dans $(2)$ : $3x + 2(185 - x) = 470 \iff x + 370 = 470 \iff \mathbf{x = 100\text{ Valencia}}$ !
    Il reste $y + z = 85$ et, dans $(3)$, $20y + 12z = 3000 - 1500 = 1500 \iff 5y + 3z = 375$. Avec $z = 85 - y$ : $5y + 3(85 - y) = 375 \iff 2y = 120 \iff \mathbf{y = 60\text{ Navel}}$ et $\mathbf{z = 25\text{ Sanguines}}$.\\
    \textbf{Étape 2 (Poids d'un sac $P$ et nombres de sacs) :} Les masses totales par variété sont $M_{\text{Val}} = 100 \times 0,15\text{ kg} = 15\text{ kg}$, $M_{\text{Nav}} = 60 \times 0,20\text{ kg} = 12\text{ kg}$ et $M_{\text{Sang}} = 25 \times 0,12\text{ kg} = 3\text{ kg}$.
    La relation sur le nombre de sacs $n_{\text{Val}} = n_{\text{Nav}} + 1$ s'écrit $\dfrac{15}{P} = \dfrac{12}{P} + 1 \iff \boxed{P = 3\text{ kg}}$, d'où $\boxed{5\text{ sacs Valencia},\; 4\text{ sacs Navel et } 1\text{ sac Sanguines}}$.

    \item \textbf{Septembre 2025, Q3 (Les filets de ballons de sport) :}
    \emph{Soient $F, B, V$ les masses (en g) d'un ballon de foot, basket, volley, et $S$ celle d'un filet vide.
    Relations entre ballons : $F = \frac{B + V}{2}$ (moyenne) et $2F + 3V = 3B$.
    On remplit $10$ filets ($a$ filets de type A contenant $3F + 2B + 3V$, et $b$ filets de type B contenant $3B + 5V$, avec $a + b = 10$) de sorte qu'il y ait autant de ballons de volley que de foot et basket réunis. Un filet A pèse $150\text{ g}$ de plus qu'un filet B, et les $10$ filets remplis pèsent $34\,750\text{ g}$.}\\
    \textbf{Résolution :}
    \begin{itemize}
        \item En remplaçant $2F = B + V$ dans $2F + 3V = 3B$, on obtient $B + 4V = 3B \iff \mathbf{B = 2V}$, puis $\mathbf{F = \frac{3}{2}V}$ ($V$ est la variable pivot !).
        \item L'écart de poids entre un filet A et un filet B donne : $(3F + 2B + 3V + S) - (3B + 5V + S) = 150 \iff 3F - B - 2V = 150$. En remplaçant en fonction de $V$ : $\frac{9}{2}V - 2V - 2V = \frac{1}{2}V = 150 \implies \boxed{V = 300\text{ g}}$, $\boxed{B = 600\text{ g}}$ et $\boxed{F = 450\text{ g}}$.
        \item Nombre de filets $a$ et $b$ ($a + b = 10$) : la condition $V_{\text{tot}} = F_{\text{tot}} + B_{\text{tot}}$ s'écrit $3a + 5b = 3a + (2a + 3b) \iff 2b = 2a \iff \mathbf{a = b = 5}$.
        \item Masse d'un filet vide $S$ : un filet A (sans filet) pèse $3(450) + 2(600) + 3(300) = 3450\text{ g}$ et un filet B pèse $3(600) + 5(300) = 3300\text{ g}$. Pour $5$ filets A et $5$ filets B : $5(3450 + 3300) + 10S = 33\,750 + 10S = 34\,750 \implies 10S = 1000 \implies \boxed{S = 100\text{ g}}$.
    \end{itemize}

    \item \textbf{Septembre 2026, Q3 (L'incendie du bois de Lauzelle : débits et rotations) :}
    \emph{$1\,060\text{ m}^3$ d'eau ont été prélevés en $149$ passages par des camions de pompiers (capacité $C$), des citernes agricoles (capacité $A$) et l'hélicoptère Charlie (capacité $H$). Le jour, les camions ont fait $3$ fois plus de passages que les citernes ($3a$ contre $a$). La nuit, seuls les camions ont fait $16$ passages. Charlie a fait $5$ passages au total. Le débit d'arrosage d'une citerne est le double de celui d'un camion ($2q$ contre $q$), mais elle se vide $3$ fois plus vite ($t/3$ contre $t$). Enfin, les $5$ passages de Charlie ont emporté autant d'eau que $8$ passages de camions.}\\
    \textbf{Résolution en 2 temps :}
    \begin{itemize}
        \item \emph{1. Nombre de passages :} $a + 3a + 16 + 5 = 149 \iff 4a = 128 \iff \mathbf{a = 32}$. Les citernes ont donc effectué $\boxed{32\text{ passages}}$ et les camions $3(32) + 16 = \boxed{112\text{ passages}}$.
        \item \emph{2. Capacités (avec $C = q\cdot t$ comme variable pivot) :}
        La capacité d'une citerne vaut $A = \text{débit} \times \text{temps} = (2q)\cdot\frac{t}{3} = \frac{2}{3}qt = \mathbf{\frac{2}{3}C}$.
        Pour l'hélicoptère : $5H = 8C \iff \mathbf{H = \frac{8}{5}C}$.
        L'équation du volume total $112C + 32A + 5H = 1060$ devient :
        \[
            112C + \frac{64}{3}C + 8C = 1060 \iff \frac{424}{3}C = 1060 \implies C = \frac{3180}{424} = \boxed{7,5\text{ m}^3},
        \]
        d'où $A = \frac{2}{3}(7,5) = \boxed{5\text{ m}^3}$ et $H = \frac{8}{5}(7,5) = \boxed{12\text{ m}^3}$.
    \end{itemize}
\end{enumerate}
\end{exemple}

% ----------------------------------------------------------------------------
\section{Problèmes ensemblistes (Diagramme de Venn), bilinéaires et suites}
% ----------------------------------------------------------------------------

\subsection{Problème de dénombrement par diagramme de Venn à 3 ensembles}

Lorsqu'un problème répartit une population de $N$ individus selon leur appartenance à trois catégories qui se chevauchent (ex. les fans photographiés avec les trois membres d'un groupe de rock), \textbf{dessinez immédiatement un diagramme de Venn à 3 cercles} et inscrivez une variable dans chacune des \textbf{7 régions disjointes} !

\begin{exemple}[Les autographes du groupe de rock \emph{The EPLighters} -- UCLouvain, Juillet 2023 Série 2, Q3]
\textbf{Résumé de l'énoncé :} Chaque fan ($N$ fans au total) repart avec exactement une photo prise avec une ou plusieurs des trois stars du groupe : Charles ($C$), Aude ($A$) et Laurent ($L$).
\begin{itemize}
    \item Charles signe au total $200$ photos, soit deux fois plus qu'Aude ($100$ photos au total pour Aude).
    \item Si Laurent avait signé $6$ fois plus de photos au total, il lui en aurait encore manqué $10$ pour couvrir tous les fans : $N = 6\,\text{Tot}_L + 10$. De plus, Laurent apparaît \textbf{seul} sur $35$ clichés.
    \item Le nombre de photos avec Charles et Aude (sans Laurent) est le double du nombre de photos avec Charles et Laurent (sans Aude). Aude et Laurent n'apparaissent \textbf{jamais} ensemble sans Charles ($0$ photo). Enfin, $5$ photos réunissent le groupe au complet ($C, A$ et $L$).
\end{itemize}
Combien y avait-il de fans $N$ à cette soirée ?

\begin{center}
\begin{tikzpicture}[scale=0.95, >=Stealth]
    % Cercle Charles
    \draw[very thick, uclblue, fill=ucllightblue, fill opacity=0.35] (-1.2, 0.7) circle (1.8);
    % Cercle Aude
    \draw[very thick, defgreen, fill=deflightgreen, fill opacity=0.35] (1.2, 0.7) circle (1.8);
    % Cercle Laurent
    \draw[very thick, exorange, fill=exlightorange, fill opacity=0.35] (0, -1.2) circle (1.8);

    \node[uclblue, font=\bfseries] at (-2.7, 2.3) {Charles ($\text{Tot}_C = 200$)};
    \node[defgreen, font=\bfseries] at (2.7, 2.3) {Aude ($\text{Tot}_A = 100$)};
    \node[exorange, font=\bfseries] at (0, -3.35) {Laurent ($\text{Tot}_L = 40 + x$)};

    % Les 7 régions disjointes
    \node[font=\bfseries] at (-1.8, 1.1) {$c = 165$};
    \node[font=\bfseries] at (1.8, 1.1) {$a = 75$};
    \node[font=\bfseries] at (0, -1.8) {$35$};
    \node[uclblue, font=\bfseries] at (0, 1.35) {$2x = 20$};
    \node[exorange, font=\bfseries] at (-0.95, -0.35) {$x = 10$};
    \node[defgreen, font=\bfseries] at (0.95, -0.35) {$0$};
    \node[warnred, font=\bfseries] at (0, 0.25) {$5$};
\end{tikzpicture}
\end{center}

\textbf{Résolution avec le diagramme de Venn :}
Notons $x = CL$ le nombre de photos avec Charles et Laurent seuls (donc $2x$ avec Charles et Aude seuls), $c$ le nombre de photos avec Charles seul et $a$ avec Aude seule. Les $7$ régions disjointes contiennent donc $c,\; a,\; 35,\; 2x,\; x,\; 0,\; 5$ :
\begin{enumerate}
    \item \textbf{Total du cercle de Charles :} $c + 2x + x + 5 = 200 \iff \mathbf{c = 195 - 3x}$.
    \item \textbf{Total du cercle d'Aude :} $a + 2x + 0 + 5 = 100 \iff \mathbf{a = 95 - 2x}$.
    \item \textbf{Total du cercle de Laurent :} $\text{Tot}_L = 35 + x + 0 + 5 = \mathbf{40 + x}$.
    \item \textbf{Nombre total de fans $N$ (somme des 7 régions disjointes) :}
    \[
        N = c + a + 35 + 2x + x + 0 + 5 = (195 - 3x) + (95 - 2x) + 3x + 40 = \mathbf{330 - 2x}.
    \]
    \item \textbf{Équation sur le total de Laurent :} $N = 6\,\text{Tot}_L + 10 = 6(40 + x) + 10 = \mathbf{250 + 6x}$.
    \item \textbf{Égalisation des deux expressions de $N$ :}
    \[
        330 - 2x = 250 + 6x \iff 8x = 80 \iff \mathbf{x = 10} \implies N = 330 - 2(10) = \boxed{310\text{ fans}}.
    \]
\end{enumerate}
\end{exemple}

\subsection{Systèmes bilinéaires (Quantité $\times$ Rendement ou Longueur $\times$ Prix unitaire)}

\begin{exemple}[Deux problèmes bilinéaires avec astuce de bloc (Septembre 2024 \& Juillet 2026)]
\begin{enumerate}
    \item \textbf{Septembre 2024, Q3 (Charles l'électricien : voie rapide vs voie algébrique) :}
    \emph{Charles utilise des rouleaux de $50\text{ m}$ de fil (rouge, bleu, noir, tous au même prix $x$\,€/m). Chaque circuit utilise deux fils de couleurs différentes de même longueur.
    Chez Aude (facture $2\,000$\,€) : uniquement des circuits rouge-bleu (longueur $y$ de rouge, donc $y$ de bleu).
    Chez Laurent (facture $6\,000$\,€) : circuits rouge-bleu et noir-bleu, avec $2$ fois plus de rouge que chez Aude ($2y$) et autant de noir que de rouge chez Aude ($y$), donc la longueur de bleu vaut $2y + y = 3y$.
    Chez Sophie (facture $14\,000$\,€) : circuits rouge-noir et bleu-noir, avec autant de bleu que chez Laurent ($3y$, donc total de fil chez Sophie $= 2z + 6y$ où $z$ est la longueur de rouge chez Sophie).
    Au total, Charles a utilisé $44$ rouleaux toutes couleurs confondues, dont $14$ rouleaux de fil rouge.}\\
    \textbf{La « voie royale » en 3 lignes :}
    \begin{itemize}
        \item Les $44$ rouleaux ($44 \times 50 = 2\,200\text{ m}$ de câble au total) ont été facturés au total $2\,000 + 6\,000 + 14\,000 = 22\,000$\,€. Le prix d'un mètre de câble est donc immédiatement :
        \[
            x = \frac{22\,000\,\text{€}}{2\,200\text{ m}} = \boxed{10\,\text{€/m}} \quad (\text{soit } 500\,\text{€ le rouleau}).
        \]
        \item Chez Aude ($2\,000$\,€, soit $200\text{ m}$ de fil moitié rouge, moitié bleu), la longueur de fil rouge vaut donc $\boxed{y = 100\text{ m}}$ ($2$ rouleaux).
        \item Laurent a deux fois plus de rouge qu'Aude, soit $200\text{ m}$ ($4$ rouleaux). Comme $14$ rouleaux rouges ($700\text{ m}$) ont été utilisés au total, il reste pour Sophie $700 - 100 - 200 = \boxed{z = 400\text{ m}}$ ($8$ rouleaux rouges) !
    \end{itemize}

    \item \textbf{Juillet 2026, Q3 (Les ruches de Charles l'apiculteur : système bilinéaire en $N \cdot p$) :}
    \emph{Soient $N$ le nombre initial de ruches, $p$ le nombre de pots produits par une ruche saine, $v$ la perte par ruche due au varroa et $f$ la perte par ruche due aux frelons asiatiques.
    Année 1 : la moitié des ruches ($\frac{N}{2}$) a le varroa $\implies$ perte totale de $180$ pots.
    Année 2 : toutes les ruches subissent les frelons $\implies$ perte d'un sixième de la production normale $Np$.
    Année 3 : $\frac{1}{4}$ des ruches est éliminé (il reste $\frac{3}{4}N$ ruches), toutes touchées par le varroa ET les frelons (production $p - v - f$ par ruche) $\implies$ récolte de $1\,755$ pots.
    Année 4 : exploitation reconstituée avec $5$ ruches de plus ($N + 5$ ruches toutes saines) $\implies$ récolte doublée par rapport à l'année 3 ($2 \times 1755 = 3\,510$ pots).}\\
    \textbf{Mise en équations et résolution par le bloc $Np$ :}
    \begin{itemize}
        \item Année 1 : $\dfrac{N}{2} v = 180 \iff \mathbf{Nv = 360}$.
        \item Année 2 : $Nf = \dfrac{1}{6}Np \iff \mathbf{f = \frac{p}{6}}$ (une ruche touchée par les frelons perd $\frac{1}{6}$ de sa production, il lui reste donc $\frac{5}{6}p$).
        \item Année 3 : $\dfrac{3}{4}N(p - v - f) = 1755 \iff \dfrac{3}{4}N\left(\frac{5}{6}p - v\right) = 1755$. En développant le produit par $N$ pour faire apparaître les blocs connus $\mathbf{Np}$ et $\mathbf{Nv = 360}$ :
        \[
            \frac{5}{8}(Np) - \frac{3}{4}(360) = 1755 \iff \frac{5}{8}(Np) = 1755 + 270 = 2025 \implies \mathbf{Np = 3240}.
        \]
        \item Année 4 : $(N + 5)p = 3510 \iff Np + 5p = 3510$. En remplaçant le bloc $Np = 3240$ :
        \[
            3240 + 5p = 3510 \iff 5p = 270 \iff \boxed{p = 54\text{ pots/ruche saine}}.
        \]
        \item On déduit aussitôt : $N = \dfrac{3240}{54} = \boxed{60\text{ ruches}}$, puis $v = \dfrac{360}{60} = 6$ (ruche avec varroa : $p - v = 54 - 6 = \boxed{48\text{ pots}}$) et $f = \dfrac{54}{6} = 9$ (ruche avec frelons : $p - f = 54 - 9 = \boxed{45\text{ pots}}$).
    \end{itemize}
\end{enumerate}
\end{exemple}

\subsection{Problèmes combinant systèmes d'équations et suites (Juillet 2022 \& Septembre 2023)}

\begin{exemple}[Deux problèmes mêlant suites et systèmes (Juillet 2022 \& Septembre 2023)]
\begin{enumerate}
    \item \textbf{Juillet 2022, Q3 (Les panneaux photovoltaïques de Charles et Laurent) :}
    \emph{Un panneau neuf produit $P = 300\text{ kWh/an}$. Les panneaux polycristallins de Charles (coût unitaire $c_p$) perdent $5\,\%$ de production par an (raison géométrique $r = 0,95$). Charles installe le plus petit nombre entier $S$ de panneaux pour que la production de la $10$\ieme{} année ($P_{10} = S \cdot 300 \cdot 0,95^9$, avec $300 \times 0,95^9 = 189,0\text{ kWh}$) couvre sa consommation $W = 2600\text{ kWh}$. Sur $10$ ans, au prix $c_e$\,€/kWh et avec un coût fixe d'installation $c_0$, Charles réalise un bénéfice net $c_C = 1560$\,€ à l'achat, ou $c_R = 300$\,€ en leasing (loyer annuel de $12\,\%$ du prix des panneaux pendant $10$ ans, soit $10 \times 0,12 = 1,2$ fois le prix d'achat des panneaux). Laurent installe $S$ panneaux monocristallins (sans dégradation, $30\,\%$ plus chers : $c_m = 1,3\,c_p$) et réalise un bénéfice net $c_L = 3390$\,€ sur $10$ ans.}\\
    \textbf{Résolution :}
    \begin{itemize}
        \item \emph{Nombre de panneaux $S$ :} À la $10$\ieme{} année, $189 \cdot S \ge 2600 \iff S \ge \frac{2600}{189} \approx 13,75 \implies \boxed{S = 14\text{ panneaux}}$.
        \item \emph{Productions sur $10$ ans :} Pour Laurent (sans dégradation), $W_L = 10 \times 14 \times 300 = \mathbf{42\,000\text{ kWh}}$. Pour Charles (suite géométrique), $W_C = 14 \times 300 \times \frac{1 - 0,95^{10}}{1 - 0,95} = \mathbf{33\,706\text{ kWh}}$.
        \item \emph{Système linéaire pour $c_p, c_m, c_e, c_0$ :}
        \[
            \begin{cases}
                c_C = W_C c_e - c_0 - S c_p = 1560 & (1) \\
                c_R = W_C c_e - c_0 - 1,2\,S c_p = 300 & (2) \\
                c_L = W_L c_e - c_0 - 1,3\,S c_p = 3390 & (3)
            \end{cases}
        \]
        En soustrayant $(1) - (2)$, $W_C c_e$ et $c_0$ s'éliminent d'un coup : $0,2\,S c_p = 1560 - 300 = 1260 \implies c_p = \frac{1260}{2,8} = \boxed{450\,\text{€}}$, d'où $c_m = 1,3 \times 450 = \boxed{585\,\text{€}}$.
        Puis $(3) - (1)$ élimine $c_0$ : $(W_L - W_C)c_e - 0,3\,S c_p = 3390 - 1560 = 1830 \iff 8294\,c_e = 1830 + 1890 = 3720 \implies c_e = \boxed{0,45\,\text{€/kWh}}$, et enfin $(1)$ donne $c_0 = \boxed{7\,257\,\text{€}}$.
    \end{itemize}

    \item \textbf{Septembre 2023, Q3 (Rénovation du living et escalier suspendu) :}
    \emph{Pièce rectangulaire de dimensions $x$ et $y$ (hauteur sous plafond $h$) avec $2$ portes ($2\text{ m} \times 1\text{ m}$ chacune) et $1$ fenêtre ($1,6\text{ m}$ de haut sur toute la longueur $y$ de la pièce).
    Parquet au sol : surface $xy = 36\text{ m}^2$.
    Plinthes au bas des 4 murs (sauf les 2 portes de $1\text{ m}$ de large) : $2x + 2y - 2 = 23 \iff 2x + 2y = 25\text{ m}$.
    Escalier central de $16$ marches régulièrement espacées : les $15$ premières marches sont suspendues chacune par $2$ barres descendant du plafond (la $16$\ieme{} est au niveau du plafond, donc $0\text{ m}$ de barre). La longueur totale des $30$ barres est de $42\text{ m}$ (donc $21\text{ m}$ pour une rangée de $15$ barres).
    Quelle quantité de peinture (en kg) faut-il pour peindre les $4$ murs (hors portes et fenêtre) en $2$ couches avec un rendement de $4\text{ m}^2/\text{kg}$ ?}\\
    \textbf{Résolution :}
    \begin{itemize}
        \item \emph{1. Hauteur sous plafond $h$ (suite arithmétique des barres d'escalier) :} La hauteur $h$ est divisée en $16$ hauteurs de marche égales à $r = \frac{h}{16}$. Les longueurs d'une rangée de $15$ barres forment la suite arithmétique $u_1 = r,\; u_2 = 2r,\; \dots,\; u_{15} = 15r$, dont la somme vaut $\frac{42}{2} = 21\text{ m}$ :
        \[
            S_{15} = \frac{15 \times 16}{2}\,r = 120\,r = 21 \implies r = \frac{21}{120} = \frac{7}{40}\text{ m} \implies h = 16\,r = \frac{14}{5} = \mathbf{2,8\text{ m}}.
        \]
        \item \emph{2. Dimensions $x$ et $y$ de la pièce (Somme-Produit !) :} On a $x + y = \frac{25}{2} = 12,5$ et $xy = 36$. Les dimensions sont les racines de $Y^2 - \frac{25}{2}Y + 36 = 0 \iff 2Y^2 - 25Y + 72 = 0$ ($\Delta = 625 - 576 = 49 = 7^2$), ce qui donne $\mathbf{y = 8\text{ m}}$ (la longueur) et $\mathbf{x = 4,5\text{ m}}$ (la largeur).
        \emph{Astuce remarquable :} pour la surface des 4 murs $2h(x + y) = 2,8 \times 25 = 70\text{ m}^2$, on avait déjà $2(x + y) = 25\text{ m}$ sans même résoudre le second degré, mais il fallait connaître la longueur $y = 8\text{ m}$ pour calculer la surface de la fenêtre !
        \item \emph{3. Surface à peindre et masse de peinture :}
        $S_{\text{murs}} = 2h(x + y) - S_{\text{fenêtre}} - 2\,S_{\text{porte}} = 2,8 \times 25 - (1,6 \times 8) - 2(2 \times 1) = 70 - 12,8 - 4 = \mathbf{53,2\text{ m}^2} = \frac{266}{5}\text{ m}^2$.
        Pour $2$ couches (soit $2 \times 53,2 = 106,4\text{ m}^2$) à $4\text{ m}^2/\text{kg}$, il faut $\dfrac{106,4}{4} = \boxed{26,6\text{ kg de peinture}}$.
    \end{itemize}
\end{enumerate}
\end{exemple}

% ----------------------------------------------------------------------------
\section{Exercices d'entraînement du Chapitre 8}
% ----------------------------------------------------------------------------

\emph{Les corrigés détaillés de tous ces exercices se trouvent au Chapitre~10.}

\begin{enumerate}[label=\textbf{Exercice 8.\arabic* :}, leftmargin=*]
    \item \textbf{(UCLouvain -- Septembre 2021, Q3 : La collecte des déchets des trois voisins)} Trois voisins, Laurent ($L$), Nadine ($N$) et Bruno ($B$), paient ensemble une facture annuelle totale de $P = 251,60$\,€ pour la collecte de leurs déchets ménagers. Laurent et Nadine ont souscrit un petit forfait annuel à $50$\,€ (utilisant un conteneur de $5\text{ kg}$, qui couvre jusqu'à $50\text{ kg}$ de déchets par an et $12$ levées par an). Bruno a souscrit un grand forfait annuel à $80$\,€ (conteneur de $10\text{ kg}$, couvrant jusqu'à $70\text{ kg}$ de déchets par an et $12$ levées par an). Tout kilogramme au-delà du forfait est facturé $0,40$\,€/kg et toute levée au-delà des $12$ levées incluses est facturée $1,20$\,€/levée.
    Sur l'année, les trois voisins ont produit au total $290\text{ kg}$ de déchets et ont totalisé $43$ levées. Chacun ne sort son conteneur que lorsqu'il est intégralement rempli, et Laurent est le seul à ne pas avoir dépassé son forfait. Déterminer les quantités annuelles de déchets $Q_L, Q_N, Q_B$ (en kg) produites par chacun des trois voisins.
    \item \textbf{(Problème d'alliages et de proportions type Question 3)} Un atelier de métallurgie de l'EPL dispose de trois lingots d'alliages de cuivre, de zinc et d'étain :
    \begin{itemize}
        \item Lingot A : $60\,\%$ de cuivre, $30\,\%$ de zinc et $10\,\%$ d'étain ;
        \item Lingot B : $50\,\%$ de cuivre, $20\,\%$ de zinc et $30\,\%$ d'étain ;
        \item Lingot C : $40\,\%$ de cuivre, $50\,\%$ de zinc et $10\,\%$ d'étain.
    \end{itemize}
    Quelle masse (en kg) faut-il prélever et refondre de chacun des trois lingots pour obtenir $100\text{ kg}$ d'un nouvel alliage contenant exactement $52\,\%$ de cuivre, $31\,\%$ de zinc et $17\,\%$ d'étain ?
    \item \textbf{(Problème de diagramme de Venn type Question 3)} Parmi les $250$ étudiants inscrits à une journée de préparation à l'EPL, chacun suit au moins l'un des trois ateliers (Algèbre, Analyse, Trigonométrie) : $140$ suivent l'Algèbre, $130$ suivent l'Analyse et $110$ suivent la Trigonométrie. On sait que $30$ étudiants suivent les trois ateliers, que le nombre d'étudiants suivant uniquement l'Algèbre et l'Analyse est le double de ceux qui suivent uniquement l'Analyse et la Trigonométrie, et que $25$ étudiants suivent uniquement l'Algèbre et la Trigonométrie. Combien d'étudiants suivent \emph{uniquement} l'atelier d'Algèbre ?
\end{enumerate}
